\section{Real roots of the Doche--Mend\`es France polynomials (\oeis{A347171})}\label{s:A347171}
The triangle \oeis{A347171} gives the coefficients of $P_n(y)=\sum_{0\le j<2^n}\mathrm{GRS}(j)\,y^{\mathrm{wt}(j)}$, where $\mathrm{GRS}$ is the Golay--Rudin--Shapiro sequence. The entry (K.~Ryde, 2021) records the conjecture of Doche and Mend\`es France that the number of real roots of $P_n$ is \oeis{A285869}$(n)$, the number of zeros of the Chebyshev polynomial $S_n$ in $(-\sqrt2,\sqrt2)$.

\paragraph{Prior work.} In their paper (Budapest 2001; version of April 2005, pp.~9--10), Doche and Mend\`es France
write that computations up to degree $200$ with PARI's \texttt{polsturm} suggest the count
$2\lceil n/4\rceil-((-1)^{n+1}+1)/2$, which equals \oeis{A285869}$(n)$, and add: ``Unfortunately we are not able to
establish this.'' We found no later proof.

\subsection{Statement}
The polynomials satisfy $P_0=1$, $P_1=1+y$ and
\begin{equation}\label{e:grsrec}
P_n(y)=(1-y)P_{n-1}(y)+2y\,P_{n-2}(y)\qquad(n\ge2),
\end{equation}
a recurrence due to Doche and Mend\`es France \cite{DMF}. Let
$A(n)=2b(n)+(n \bmod 2)$ with $b(n)=\lfloor n/2\rfloor-\lfloor (n+1)/4\rfloor$; this is the closed form of
\oeis{A285869}$(n)$. Put $L=-3-2\sqrt2$ and $R=-3+2\sqrt2$.

\begin{theorem}\label{t:A347171}
For every $n\ge0$ the real roots of $P_n$ are simple and their number is $A(n)$. More precisely, for
$n\ge2$, with $q=(n+3)/4$ and $K=\lceil q\rceil-1$, the number of roots of $P_n$ is
\[
\begin{array}{ll}
0 & \text{in } (-\infty,L]\ \text{and in}\ [R,0],\\
K-1+[q\in\mathbb Z] & \text{in } (L,-1),\\
{}[q\in\mathbb Z] & \text{at } y=-1,\\
K & \text{in } (-1,R),\\
{}[n \text{ even}] & \text{in } (0,\infty).
\end{array}
\]
\end{theorem}
\noindent Here $[\,\cdot\,]$ is $1$ if the condition holds and $0$ otherwise. The cases $n=0,1$ are trivial
($P_0=1$, $P_1=1+y$).

\subsection{The Chebyshev form for \texorpdfstring{$y<0$}{y<0}}
Let $S_n$ be defined by $S_{-1}=0$, $S_0=1$, $S_n(x)=xS_{n-1}(x)-S_{n-2}(x)$, so that
$S_n(2\cos\theta)=\sin((n+1)\theta)/\sin\theta$ and $S_n(2\cosh\eta)=\sinh((n+1)\eta)/\sinh\eta$.

\begin{lemma}\label{l:cheb}
For $y<0$ put $t=\sqrt{-2y}>0$ and $x=(1-y)/t=t/2+1/t$. Then $P_n(y)=t^n\bigl(S_n(x)-t\,S_{n-1}(x)\bigr)$.
\end{lemma}
\begin{proof}
Both sides equal $1$ for $n=0$ and $tx-t^2=1+y$ for $n=1$. Writing $P_n=t^nQ_n$, the recurrence \eqref{e:grsrec}
becomes $Q_n=xQ_{n-1}-Q_{n-2}$ because $tx=1-y$ and $t^2=-2y$.
\end{proof}
Since $x=t/2+1/t\ge\sqrt2$, we have $x<2$ exactly when $t\in(2-\sqrt2,2+\sqrt2)$, that is, $y\in(L,R)$.

\begin{lemma}\label{l:outside}
For $n\ge1$, $P_n$ has no root in $(-\infty,L]\cup[R,0]$.
\end{lemma}
\begin{proof}
$P_n(0)=1$. At $y=L$ or $R$ we have $x=2$, $S_n(2)=n+1$ and $P_n=t^n(n+1-tn)\ne0$ since $t=2\pm\sqrt2$
is irrational. Otherwise $x>2$; write $x=2\cosh\eta$, $\eta>0$. By Lemma~\ref{l:cheb}, $P_n(y)=0$ if and only if
$\rho:=\sinh((n+1)\eta)/\sinh(n\eta)=t$. Now $\rho=\cosh\eta+\coth(n\eta)\sinh\eta$ and
$1<\coth(n\eta)\le\coth\eta$, so $e^\eta<\rho\le2\cosh\eta=x$. If $t>2+\sqrt2$ then $\rho\le x=t/2+1/t<t$.
If $t<2-\sqrt2$ then $\rho>e^\eta>1>t$.
\end{proof}

\subsection{The phase on \texorpdfstring{$(L,R)$}{(L,R)}}
For $y\in(L,R)$ write $x=2\cos\theta$ with $\theta\in(0,\pi/4]$, and let $c=\cos\theta$, $s=\sin\theta$,
$w=\sqrt{2\cos2\theta}=\sqrt{4c^2-2}$. Since $t^2-2xt+2=0$, we get $t=2c\pm w$:
\begin{itemize}
\item branch 1: $t_1=2c+w\in(\sqrt2,2+\sqrt2)$, which describes $y\in(L,-1)$;
\item branch 2: $t_2=2c-w\in(2-\sqrt2,\sqrt2)$, which describes $y\in(-1,R)$.
\end{itemize}
On each branch $\theta\mapsto y$ is a bijection from $(0,\pi/4)$; $\theta=\pi/4$ is $y=-1$, and
$\theta\to0$ is $y\to L$ on branch 1 and $y\to R$ on branch 2. Since
$\sin((n+1)\theta)-t\sin(n\theta)=\operatorname{Im}\bigl(e^{in\theta}(e^{i\theta}-t)\bigr)$, Lemma~\ref{l:cheb} gives
\begin{equation}\label{e:phase}
P_n(y)=t^n\,\bigl|e^{i\theta}-t\bigr|\,\frac{\sin\Phi(\theta)}{\sin\theta},\qquad
\Phi(\theta)=n\theta+\beta(\theta),\quad \beta(\theta)=\arg\bigl(e^{i\theta}-t\bigr)\in(0,\pi),
\end{equation}
and $|e^{i\theta}-t|^2=1-2tc+t^2=t^2/2>0$. So the roots in $(L,R)$ are the solutions of $\Phi\in\pi\mathbb Z$.

\begin{lemma}\label{l:b2}
On branch 2, $\beta_2$ is strictly increasing from $0$ to $3\pi/4$. Hence $\Phi_2=n\theta+\beta_2$ is
strictly increasing from $0$ to $q\pi$.
\end{lemma}
\begin{proof}
$\cot\beta_2=(c-t_2)/s=(w-c)/s$. With $w'=-4sc/w$, the derivative of $(w-c)/s$ is $N/s^2$ where
$N=(w'+s)s-(w-c)c$, and $Nw=-4s^2c+w-(4c^2-2)c=w-2c<0$ because $w^2=4c^2-2<4c^2$. So $\cot\beta_2$ is
strictly decreasing. At $\theta\to0$, $\cot\beta_2\to+\infty$; at $\theta=\pi/4$, $\cot\beta_2=-1$.
\end{proof}

\begin{lemma}\label{l:b1}
On branch 1, $\beta_1'(\theta)=1-2c/w$. Hence $\beta_1$ and $\Phi_1=n\theta+\beta_1$ are strictly concave,
$\Phi_1(0^+)=\pi$, $\Phi_1(\pi/4)=q\pi$, $\Phi_1'(\theta)\to-\infty$ as $\theta\to\pi/4$, and
$\Phi_1<q\pi+\pi/4$ on $(0,\pi/4)$.
\end{lemma}
\begin{proof}
$e^{i\theta}-t_1=-(c+w)+is$ lies in the second quadrant, so $\beta_1=\pi-\arctan\bigl(s/(c+w)\bigr)$.
Differentiating,
\[
\frac{d}{d\theta}\arctan\frac{s}{c+w}=\frac{1+cw-sw'}{(c+w)^2+s^2}=\frac{w+2c}{w\,(4c^2-1+2cw)} ,
\]
using $-sw'=4s^2c/w$ and $w+cw^2+4s^2c=w+2c$. Since $(2c+w)^2=2(4c^2-1+2cw)$ and $(2c+w)(2c-w)=2$, this
equals $2/\bigl(w(2c+w)\bigr)=(2c-w)/w=2c/w-1$. Now $2c/w=2/\sqrt{4-2/c^2}$ increases strictly with $\theta$,
so $\beta_1'$ decreases strictly. The boundary values are $\beta_1(0^+)=\pi$ and $\beta_1(\pi/4)=3\pi/4$, and
$\beta_1'\to-\infty$ as $w\to0$. Finally $\theta<\pi/4$ and $\beta_1<\pi$ give $\Phi_1<n\pi/4+\pi=q\pi+\pi/4$.
\end{proof}

\subsection{Positive \texorpdfstring{$y$}{y}}
\begin{lemma}\label{l:pos}
For $y>0$, $P_n$ has exactly one root, which is simple, if $n\ge2$ is even, and no root if $n$ is odd.
\end{lemma}
\begin{proof}
Let $\rho_1>0>\rho_2$ be the roots of $f(r)=r^2-(1-y)r-2y$. Since $f(1)=-y<0<2=f(2)$, $1<\rho_1<2$. The
recurrence \eqref{e:grsrec} at $n=1$ gives $P_{-1}=1$, hence
$P_n=\bigl(\rho_1^{n+1}(\rho_2-1)+\rho_2^{n+1}(1-\rho_1)\bigr)/(\rho_2-\rho_1)$, and $P_n(y)=0$ if and only if
$\rho_1^{n+1}(1-\rho_2)=\rho_2^{n+1}(1-\rho_1)$. The left side is positive. If $n$ is odd the right side is
negative. If $n$ is even, parametrize by $\rho=\rho_1\in(1,2)$: then $y=\rho(\rho-1)/(2-\rho)$ is an increasing
bijection onto $(0,\infty)$, $|\rho_2|=2(\rho-1)/(2-\rho)$, and the equation becomes
$\kappa^{n+2}=2/(2-\rho)^2$ with $\kappa=2(\rho-1)/(\rho(2-\rho))$. Let
$g(\rho)=(n+2)\log\kappa+2\log(2-\rho)-\log2$. Then
$g'(\rho)=(n+2)\bigl(\tfrac1{\rho-1}-\tfrac1\rho\bigr)+\tfrac{n}{2-\rho}>0$, $g\to-\infty$ at $1^+$ and, for
$n\ge1$, $g\to+\infty$ at $2^-$. So there is exactly one root, and it is simple.
\end{proof}

\subsection{Proof of Theorem \ref{t:A347171}}
Let $n\ge2$, so $q>1$. By \eqref{e:phase} and Lemma~\ref{l:b2}, the roots in $(-1,R)$ correspond to the integers
$k$ with $0<k<q$: there are $K$ of them. At $y=-1$, $\Phi=q\pi$, so $-1$ is a root if and only if
$q\in\mathbb Z$, that is, $n\equiv1\pmod4$. On branch 1, by Lemma~\ref{l:b1}, $\Phi_1$ is strictly concave with
boundary values $\pi<q\pi$, its maximum $M$ is interior with $q\pi<M<q\pi+\pi/4$, and the only multiple of
$\pi$ in $[q\pi,M]$ is $q\pi$ itself when $q\in\mathbb Z$ (otherwise the next multiple is at least
$(q+1/4)\pi$). By concavity, each $k\pi$ with $1<k<q$ is attained exactly once, on the increasing part; if
$q\in\mathbb Z$, $q\pi$ is attained once in the interior and once at the endpoint $\theta=\pi/4$; and $\pi$ is
not attained in the interior. This gives $K-1+[q\in\mathbb Z]$ roots in $(L,-1)$. All these roots are simple,
because $\Phi'\ne0$ there (at $y=-1$ one uses the parameter $t$, in which $\theta'(\sqrt2)=0$ and $\beta'\ne0$).
With Lemmas~\ref{l:outside} and~\ref{l:pos}, the total is $2K-1+2[q\in\mathbb Z]+[n\text{ even}]$.
Writing $n=4m+r$: for $r=0,1,2,3$ we get $(q,K)=(m+\tfrac34,m)$, $(m+1,m)$, $(m+\tfrac54,m+1)$,
$(m+\tfrac32,m+1)$ and totals $2m$, $2m+1$, $2m+2$, $2m+1$, while $b(n)=m,m,m+1,m$. In each case the total
is $A(n)$. \qed


\begin{remark}
Every formula above was checked numerically, and the table of Theorem~\ref{t:A347171} was confirmed for
$1\le n\le300$ with exact Sturm sequences, the signs at $L$ and $R$ being computed in $\mathbb Q(\sqrt2)$.
The scripts are in the repository \url{https://github.com/joaquimmj/oeis-walnut}, folder
\texttt{conjectures/A347171}.
\end{remark}



\section{The Josephus problem: Hadjicostas' conjectures}\label{s:A005427}
\subsection{The general conjecture}
Put $m$ people in a circle, numbered $1,\dots,m$, and remove every $k$-th person, starting the count at $1$.
Let $J_k(m)$ be the survivor; it satisfies $J(1)=1$ and $J(m)=((J(m-1)+k-1)\bmod m)+1$. Let $T(1)=1$ and
$T(n)=\lceil(T(1)+\dots+T(n-1))/(k-1)\rceil$, and $S(n)=T(1)+\dots+T(n)$. In a comment of July 22, 2020 in
\oeis{A005427}, P.~Hadjicostas conjectured that the $m\ge k$ with $J_k(m)\in\{1,\dots,k-1\}$ are exactly the
values $T(n)\ge k$, and wrote that he did not have a general proof. The cases in the OEIS are $k=2$ (powers
of $2$), $k=3$ (\oeis{A073941}, \oeis{A005428}), $k=4$ (\oeis{A072493}, \oeis{A005427}) and $k=5,\dots,9$
(\oeis{A120160}, \oeis{A120170}, \oeis{A120178}, \oeis{A120186}, \oeis{A120194}).

\begin{theorem}\label{t:josephus} The conjecture is true for every $k\ge2$. Moreover, for $T(n)\ge k$,
\[ J_k(T(n))=r_n+1,\qquad r_n=(k-1)T(n)-S(n-1)\in\{0,\dots,k-2\}. \]\end{theorem}

\begin{proof} The number $r_n$ is the remainder of the ceiling, so $T(n+1)=T(n)+t_n$ with
$t_n=\lceil(T(n)-r_n)/(k-1)\rceil$, and $r_{n+1}=r_n+(k-1)t_n-T(n)$.

\emph{Small values and the base case.} While $T(n)\le k-1$, we have $t_n\in\{0,1\}$, since
$T(n)-r_n\le k-1$; so $T$ takes each of the values $1,\dots,k-1$, possibly several times. For $v\le k-1$ let
$Q(v)$ be the partial sum $S$ at the last occurrence of $v$. The next term is $v+1$ exactly when the sum
exceeds $(k-1)v$, and the steps are equal to $v$, so $Q(v)\equiv Q(v-1)\pmod v$ and $Q(v)$ is the least
integer $>(k-1)v$ in this class; hence $(k-1)v<Q(v)\le kv$, with $Q(0)=0$ and $Q(1)=k$.

\emph{Claim:} $J(v)=kv+1-Q(v)$ for $1\le v\le k-1$. For $v=1$ both sides are $1$. If it holds for $v-1$, then
$J(v-1)+k-1=kv-Q(v-1)$, and since $Q(v)\equiv Q(v-1)\pmod v$ and $0\le kv-Q(v)\le v-1$,
$(J(v-1)+k-1)\bmod v=kv-Q(v)$, which gives $J(v)=kv-Q(v)+1$.

After the last occurrence of $k-1$ the next term is $T(n_0)=k$, because $(k-1)^2<Q(k-1)\le k(k-1)$. Then
$r_{n_0}=k(k-1)-Q(k-1)\in[0,k-2]$, and by the claim $J(k-1)+k-1=k(k-1)-Q(k-1)=r_{n_0}<k$, so
$J(k)=r_{n_0}+1\le k-1$. Thus $T(n_0)=k$ is the first value $\ge k$, it is in the Josephus set, and the formula
holds there.

\emph{Induction.} For $m\ge k-1$ we have $J(m-1)\le m-1$, so $J(m-1)+k-1<2m$ and the count wraps around at
most once per step. Suppose $J(m_0)=j=r_n+1$ with $m_0=T(n)\ge k$. For $s\le t_n-1=\lfloor(m_0-j)/(k-1)\rfloor$
we have $j+(k-1)s\le m_0$ and $J(m_0+s)=j+ks\ge k+1$: none of these $m$ is in the set. For $s=t_n$ we have
$j+kt_n-1\ge m_0+t_n$, so the count wraps and
\[ J(m_0+t_n)=j+kt_n-(m_0+t_n)=r_n+1+(k-1)t_n-T(n)=r_{n+1}+1\in[1,k-1]. \]
Since $m_0+t_n=T(n+1)$ and $t_n\ge1$ (because $T(n)\ge k>r_n$), the induction goes on, and the $m\ge k$ with
$J(m)\le k-1$ are exactly the $T(n)\ge k$. \end{proof}

\begin{remark}[Prior work] The conjecture is in fact a short corollary of the classical formula of Odlyzko
and Wilf~\cite{OW}, also in Graham--Knuth--Patashnik's \emph{Concrete Mathematics} (\S3.3):
$J_k(m)=km+1-D_j$, where $D_0=1$, $D_j=\lceil kD_{j-1}/(k-1)\rceil$ and $j$ is least with $D_j>(k-1)m$
(we checked it for $k\le8$, $m<400$). Indeed $S(n)=D_{n-1}$, so $T(n+1)=D_n-D_{n-1}=\lceil D_{n-1}/(k-1)\rceil$. If
$m=\lceil D_{j-1}/(k-1)\rceil\ge k$, then $D_j=D_{j-1}+m$ is the least $D$ above $(k-1)m$ and
$J_k(m)=(k-1)m+1-D_{j-1}\in[1,k-1]$. Conversely, if $J_k(m)\le k-1$ then $D_j\ge k(m-1)+2$; with
$c=\lceil D_{j-1}/(k-1)\rceil\le m$, the assumption $c\le m-1$ would give both $D_{j-1}\ge(k-1)m-k+3$ and
$D_{j-1}\le(k-1)(m-1)$, a contradiction, so $m=c$ is a value of $T$. The proof above is independent of that formula.
\end{remark}
\begin{remark} Hadjicostas expected Schuh's 1968 method for $k=3$ to help. The argument above is direct: the
Josephus recursion, observed at the times when the survivor is small, coincides with the ceiling recursion
defining $T$. This is the principle of the Burde--Thériault algorithm, with the explicit identification
$j=r+1$, which closes the proof. We have not checked whether Burde (1987), Thériault or Odlyzko--Wilf
(1991)~\cite{OW} contain the general statement. The script \texttt{A005427/check.py} checks the equality of the
sets for $k=2,\dots,30$ and $m\le10^6$, the claim for $k\le60$, and the formula for $J$ for $T(n)\le2000$,
$k\le60$.\end{remark}

\subsection{Base \texorpdfstring{$(k+1)/k$}{(k+1)/k} digit counts}\label{s:A120160}
Let $k\ge1$, $p=k+1$, $T(1)=1$ and $T(n)=\lceil S(n-1)/k\rceil$. For \oeis{A120160}, $k=4$ and the base is
$5/4$. Hadjicostas (2020) conjectured in \oeis{A120160}: (Josephus) from $a(11)=5$ on, the terms are the
$m$ whose Josephus survivor for step $5$ is in $\{1,2,3,4\}$; (1) $a(n)$ is the number of multiples of $k$
with $n-1$ digits in base $p/k$ (in the sense of \oeis{A024634}); (2) $a(n)$ is the number of integers $\ge0$
with $n$ digits, divided by $p$ ($0$ has one digit). In base $p/k$ the last digit of $N$ is $d=N\bmod p$ and the
other digits are those of $N_1=k(N-d)/p$; for instance $36=4321031$ in base $5/4$.

\begin{theorem}\label{t:A120160} For every $k\ge1$: the Josephus conjecture holds (it is the case $k+1$ of
Theorem~\ref{t:josephus}); conjecture (2) is true for every $n\ge1$; conjecture (1) is true if only multiples
$>0$ are counted, and with $0$ counted it fails only for $n=2$.\end{theorem}
\begin{proof} For $N\ge1$ the parent of $N$ is $P(N)=k\lfloor N/p\rfloor$, and the length satisfies
$\ell(N)=1+\ell(P(N))$, $\ell(0)=0$. Since $P$ is nondecreasing, so is $\ell$, and the integers of length $n$
form an interval $I_n=[L_n,L_{n+1})$, with $I_1=\{1,\dots,k\}$. A multiple $kM\ge1$ has exactly $p$ children
$pM,\dots,pM+k$, and every $N\ge p$ has a parent $\ge k$. So $I_{n+1}$ is the union of the blocks $[pM,pM+k]$
for $kM\in I_n$. Hence $c_{n+1}=p\,m_n$, where $c_n=|I_n|$ and $m_n$ is the number of multiples of $k$ in $I_n$,
and $L_{n+1}=pU_n$ with $U_n=\lceil L_n/k\rceil$. Then $U_{n+1}=\lceil pU_n/k\rceil=U_n+\lceil U_n/k\rceil$ and
$m_n=U_{n+1}-U_n=\lceil U_n/k\rceil$. With $U_1=1=S(1)$ this is the recursion $S(n+1)=S(n)+\lceil S(n)/k\rceil$,
so $U_n=S(n)$ and $m_n=T(n+1)$, which is conjecture (1) for positive multiples. The number of integers
$\ge0$ with $n$ digits is $c_1+1=p=pT(1)$ for $n=1$ and $c_n=pm_{n-1}=pT(n)$ for $n\ge2$, which is
conjecture (2). \end{proof}
For $k=3$ (\oeis{A072493}) J.~Song announced a proof of the two digit conjectures in a comment of October
14, 2026; the proof above holds for every $k$. Check: \texttt{A120160/base\_check.py} ($k=2,\dots,12$, all
lengths below $16$ to $35$).

\section{Schreier sets and Kotesovec's exponent (\oeis{A371176})}\label{s:A371176}
Let $a(n)$ be the $n$-th integer $k\ge1$ with $\oeis{A000120}(k)\le\oeis{A001511}(k)$, that is, whose binary
weight is at most the $2$-adic valuation plus one. Kotesovec (2025) conjectured that
$\limsup\log a(n)/\log n=\log2/\log\varphi$.
\begin{theorem}\label{t:A371176} $\log a(n)/\log n\to\log2/\log\varphi$; in particular the limit exists.\end{theorem}
\begin{proof} (Branicky, in the entry.) To $k$ associate the set $X$ of positions of its $1$ bits, the least
significant having position $1$. Then $\oeis{A000120}(k)=|X|$ and $\oeis{A001511}(k)=\min X$, so $k$ is a term
if and only if $X$ is a non-empty Schreier set ($|X|\le\min X$), and $2^{m-1}\le k<2^m$ if and only if
$\max X=m$.

Let $s(m)$ be the number of Schreier sets with maximum $m$. We show $s(m)=F_m$. Clearly $s(1)=s(2)=1$. For
$m\ge3$, the shift $X\mapsto X-1$ is a bijection from the Schreier sets of maximum $m$ with $|X|<\min X$ onto
all Schreier sets of maximum $m-1$ (the inverse raises the minimum, so it preserves the condition). The
remaining sets have $|X|=\min X=j$: they consist of $j$, $m$, and $j-2$ elements of $(j,m)$, so there are
$\sum_j\binom{m-j-1}{j-2}=F_{m-2}$ of them, by $\sum_i\binom{N-i}{i}=F_{N+1}$. Hence $s(m)=s(m-1)+F_{m-2}$ and
$s(m)=F_m$ by induction.

So $N(m):=\#\{n:a(n)<2^m\}=\sum_{j\le m}F_j=F_{m+2}-1$. If $a(n)$ has $m$ bits, then $N(m-1)<n\le N(m)$, so
$\varphi^m/C\le n\le C\varphi^m$ and $2^{m-1}\le a(n)<2^m$ for a constant $C$. Hence $\log a(n)=m\log2+O(1)$,
$\log n=m\log\varphi+O(1)$, and the ratio tends to $\log2/\log\varphi$. \end{proof}
The count $s(m)=F_m$ of Schreier sets with maximum $m$ is a known fact (A.~Bird); the contribution is to
connect it with Kotesovec's conjecture. Check: \texttt{A371176/check.py} (each step for $m\le22$). The convergence is slow (the ratio is $1.419$ at
$n=46\,367$) because of the $O(1)/m$ term.

\section{Nested recursions of Gutkovskiy}\label{s:A366860}
For $k\ge1$ let $a_k(0)=0$ and $a_k(n)=n-a_k(\lfloor a_k(n-1)/k\rfloor)$. Gutkovskiy (2023) conjectured
$a_k(n)\sim c_k n$ for $k=3,4,5$ (\oeis{A366860}, \oeis{A366870}, \oeis{A366871}). For $k=2$ this is
\oeis{A286389}, settled by Shallit with \Walnut, and $k=1$ is Hofstadter's $G$-sequence \oeis{A005206}.
\begin{theorem}\label{t:A366860} Let $\sigma_k:0\mapsto1^k$, $1\mapsto1^k0$, and let $w=w_1w_2\cdots$ be its
fixed point beginning with $1$. Then $a_k(n)=w_1+\dots+w_n$ for all $n\ge0$. In particular
$a_k(n)/n\to c_k=(\sqrt{k^2+4k}-k)/2$, the positive root of $c^2+kc-k$: $c_3=(\sqrt{21}-3)/2$,
$c_4=2(\sqrt2-1)$, $c_5=(3\sqrt5-5)/2$.\end{theorem}
\begin{proof} Let $b(n)=w_1+\dots+w_n$. Since the recursion determines $a_k$, it suffices to show
$b(n)=n-b(\lfloor b(n-1)/k\rfloor)$ for $n\ge1$. As $w=\sigma(w)$, $w$ is the concatenation of the blocks
$\sigma(w_1)\sigma(w_2)\cdots$; each block begins with $k$ letters $1$ and, when $w_j=1$, ends with a $0$, its only
$0$. Fix $n\ge1$ and write the prefix of length $n-1$ as $j$ complete blocks followed by $r$ letters of block
$j+1$, $0\le r<|\sigma(w_{j+1})|$. Then $b(n-1)=kj+r$ with $r\le k$, and $n-b(n)$, the number of $0$s in the
prefix of length $n$, is the number of complete blocks of letter $1$ that it contains. If $r<k$, then
$\lfloor b(n-1)/k\rfloor=j$ and letter $n$ is a $1$ of block $j+1$; that block is complete only if it is
$\sigma(0)$, which has no $0$; so $n-b(n)=\#\{i\le j:w_i=1\}=b(j)$. If $r=k$, block $j+1$ is $\sigma(1)$ and letter
$n$ is its final $0$; then $\lfloor b(n-1)/k\rfloor=j+1$ and $n-b(n)=b(j+1)$. In both cases
$n-b(n)=b(\lfloor b(n-1)/k\rfloor)$.

The matrix of $\sigma_k$ is primitive, so the frequency of $1$ exists and equals $\lambda/(1+\lambda)$ with
$\lambda=(k+\sqrt{k^2+4k})/2$ (Lemma~\ref{l:perron}), which gives $c_k$. \end{proof}
\begin{corollary} For every $n\ge1$, the partial sums of row $n$ of \oeis{A368281} (the fixed point of
$0\mapsto1^n$, $1\mapsto1^n0$) form row $n$ of \oeis{A368282} ($b(0)=0$, $b(k)=k-b(\lfloor b(k-1)/n\rfloor)$), as
conjectured in 2024.\end{corollary}
The second eigenvalue $-k/\lambda$ has modulus less than~$1$, so $a_k(n)-c_kn$ is even bounded. For $k=2$
this gives again Irvine's conjecture proved by Shallit. Check: \texttt{A366860/guess.py} ($k=2,\dots,7$ on
$10^6$ terms, $k\le30$ on $10^5$ terms) and \texttt{A368281/check.py}.

\section{Zeros of \texorpdfstring{$a(3n)$}{a(3n)} in Stern-like sequence \oeis{A005590}}\label{s:A005590}
Let $a(0)=0$, $a(1)=1$, $a(2n)=a(n)$, $a(2n+1)=a(n+1)-a(n)$. Stephan (2003) conjectured that $a(3n)=0$ if and
only if $n$ is fibbinary (\oeis{A003714}: no two adjacent $1$ bits).
\begin{theorem}\label{t:A005590} This is true. Moreover, if $n$ is not fibbinary, $a(3n)$ is even and
$|a(3n)|\ge2$.\end{theorem}
\begin{proof} Let $u(n)=(a(3n),a(3n+1),a(3n+2),a(3n+3))$. The definition gives
$u(2n)=B_0u(n)=(a_0,a_1-a_0,a_1,a_2-a_1)$ and $u(2n+1)=B_1u(n)=(a_2-a_1,a_2,a_3-a_2,a_3)$, with $u(0)=(0,1,1,0)$;
so $u(n)$ is computed by reading the bits of $n$ from the most significant one. Let $P_m=(0,1,1,-2m)$ and
$Q_m=(0,1,-2m-1,-2m)$. Then $B_0P_m=P_0$, $B_1P_m=Q_m$, $B_0Q_m=P_{m+1}$. Since $u(0)=P_0$, as long as the bits read
contain no $11$, the state is some $P_m$ (after a $0$) or $Q_m$ (after a $1$); in particular $a(3n)=0$ if $n$ is
fibbinary.

Let $I$ be the set of vectors with $C_1=a_0(a_0-a_1)>0$, $C_2=a_1(a_1-a_2)>0$, $C_3=a_2(a_2-a_3)>0$. Under $B_0$
the three quantities become $C_1+a_0^2$, $C_1$ and $C_2+a_1^2$; under $B_1$ they become $C_2$, $C_3+a_2^2$ and
$C_3$. So $I$ is stable under $B_0$ and $B_1$. At the first $11$, the state is $B_1Q_m=(-2m-2,-2m-1,1,-2m)$,
which is in $I$ ($C_1=2m+2$, $C_2=(2m+1)(2m+2)$, $C_3=2m+1$). Hence the state stays in $I$, and $I$ forces
$a_0\ne0$: $a(3n)\ne0$ if $n$ is not fibbinary. Finally the parity pattern (even, odd, odd, even) is stable
under $B_0$ and $B_1$ and holds for $(0,1,1,0)$, so $a(3n)$ is even, and $|a(3n)|\ge2$ in the second case.
\end{proof}
Checks: \texttt{A005590/check\_symbolic.py} (the identities, with sympy) and \texttt{A005590/check.py}
(the statement and membership in $I$ for all $n<2^{21}$).

\section{Golomb's sequence: Song's lower bound}\label{s:A001462}
Let $a$ be Golomb's sequence ($n$ appears $a(n)$ times) and $\alpha=\varphi-1=1/\varphi$. Song (2021)
conjectured $a(n)\ge n^{\varphi-1}$ for all $n\ge1$.
\begin{theorem}\label{t:A001462} Let $S(m)=a(1)+\dots+a(m)=\oeis{A001463}(m)$. For all $m\ge1$,
$m^\varphi/\varphi\le S(m)\le m^\varphi$. Consequently $a(n)\ge n^{\varphi-1}$ for all $n\ge1$, with equality only
at $n=1$.\end{theorem}
\begin{proof} $S(m)$ is the last position of the value $m$, so $a(n)=m$ for $S(m-1)<n\le S(m)$, and since
$\varphi\alpha=1$, $a(n)\ge n^\alpha$ for all $n$ is equivalent to $S(m)\le m^\varphi$ for all $m$. We prove both
inequalities by strong induction. They are checked by computation for $m\le20\,000$ (60 digits). Let $m>30$
and assume the bounds below $m$. For $j\le m$ put $v=a(j)$; then $v\le a(m)<m$.

Upper bound: $j>S(v-1)\ge(v-1)^\varphi/\varphi$ gives $v<1+(\varphi j)^\alpha$, hence
$S(m)<m+\varphi^\alpha\sum_{j\le m}j^\alpha\le m+\varphi^\alpha\int_1^{m+1}x^\alpha\,dx<m+\varphi^{\alpha-1}(m+1)^\varphi$.
This is at most $m^\varphi$ if and only if $h(m)=\varphi^{\alpha-1}(1+1/m)^\varphi+m^{1-\varphi}\le1$; both terms of
$h$ decrease, and $h(30)=0.99964\ldots<1$.

Lower bound: $j\le S(v)\le v^\varphi$ gives $a(j)=v\ge j^\alpha$, so
$S(m)\ge\sum_{j\le m}j^\alpha\ge\int_0^mx^\alpha\,dx=m^\varphi/\varphi$. \end{proof}
The asymptotic constant $S(m)/m^\varphi\to\varphi^{-\varphi(2-\varphi)}\approx0.743$ (Fine, Vardi) leaves a margin;
the argument does not use it. The other conjecture of the entry (van Lamoen, 2005) follows from relation (4)
of Claes--Miyamoto~\cite{CM} and is not counted.

\section{Wythoff rows and Knuth's Fibonacci product (\oeis{A357097})}\label{s:A357097}
The entry defines a multiplication $A(\cdot,\cdot)$ of the rows of the extended Wythoff array; let
$\oeis{A019586}(n)$ be the row (from $0$) containing $n$, and $x\circ y=\oeis{A101330}(x,y)$ Knuth's Fibonacci
product: if $x=\sum F(c_i)$ and $y=\sum F(d_j)$ are Zeckendorf representations (indices $\ge2$), then
$x\circ y=\sum F(c_i+d_j)$. The conjecture is $A(\oeis{A019586}(x),\oeis{A019586}(y))=\oeis{A019586}(x\circ y)$.
\begin{theorem} The conjecture is true.\end{theorem}
\begin{proof} Let $\omega=1/\varphi$, so $\omega^2=1-\omega$. The law of the entry,
$(j_1,k_1)\circ(j_2,k_2)=(j_1j_2+k_1k_2,\ j_1k_2+k_1j_2-k_1k_2)$, is the multiplication of $j+k\omega$ in
$\mathbb{Z}[\varphi]$. Since $(j+k\omega)\varphi=(j+k)+j\omega$, moving from the pair $(eW(n,m+1),eW(n,m))$ to the next
pair multiplies by $\varphi$. So row $n$ corresponds to $S_n=\{\beta_n\varphi^m:m\in\mathbb{Z}\}$ for some
$\beta_n\in\mathbb{Z}[\varphi]$, and the $S_n$ are pairwise disjoint (two consecutive terms determine the row).

For $n\ge1$ with Zeckendorf representation $n=\sum F(c_i)$, let $\theta(n)=\sum\varphi^{c_i}$. As
$\varphi^c=F(c-1)+F(c)\varphi$, $\theta(n)=a+n\varphi$ with $a=\sum F(c_i-1)$. In the Wythoff array the successor of
$n$ in its row is $s(n)=\sum F(c_i+1)$ (the Zeckendorf array is the Wythoff array, Kimberling 1995), the pair
$(s(n),n)$ lies in row $\oeis{A019586}(n)$, and its element is $s(n)+n\omega=(s(n)-n)+n\varphi=\theta(n)$.

$\theta$ is multiplicative for $\circ$: the coefficient of $\varphi$ in $\theta(x)\theta(y)=\sum\varphi^{c_i+d_j}$ is
$\sum F(c_i+d_j)=x\circ y$. For the integer part, let $\psi=-1/\varphi$. The conjugate of $\theta(n)$ is
$\sum\psi^{c_i}$, a finite sum over non-adjacent exponents $\ge2$, hence strictly between
$\psi^3+\psi^5+\dots=-1/\varphi^2$ and $\psi^2+\psi^4+\dots=1/\varphi$. The interval $I=(-1/\varphi^2,1/\varphi)$ has length
$1$, so for given $N$ there is at most one integer $a$ with $a+N\psi\in I$, and $\theta(N)$ is that element. The
conjugate of $\theta(x)\theta(y)$ is a product of two numbers of $I$, hence in $(-1/\varphi^3,1/\varphi^2)\subset I$.
So $\theta(x)\theta(y)=\theta(x\circ y)$.

Let $u=\oeis{A019586}(x)$ and $v=\oeis{A019586}(y)$. Then $\theta(x)\in S_u$, $\theta(y)\in S_v$, the product lies in
$S_{A(u,v)}$ by definition of $A$, and it equals $\theta(x\circ y)\in S_{\oeis{A019586}(x\circ y)}$. The rows being
disjoint, $A(u,v)=\oeis{A019586}(x\circ y)$. \end{proof}
The multiplicativity of $\theta$ is Arnoux's remark (1989) on the Fibonacci product, cited in the entry; the
argument above is self-contained. Check: \texttt{A357097/check.py} ($1\le x\le y<400$).

\section{Integers with \texorpdfstring{$n$}{n} prime factors (\oeis{A078841})}\label{s:A078841}
Let $a(n)$ be the $(n+1)$-th integer $m$ with $\Omega(m)=n$ (prime factors counted with multiplicity):
$1,3,9,20,54,112,\dots$ R.~G.~Wilson~v (2007) conjectured $a(n+1)/a(n)\to2$.
\begin{theorem} Let $V=\{u/2^{\Omega(u)}:u\text{ odd}\ge1\}$, with multiplicity, in increasing order
$t_1\le t_2\le\cdots$. Then $a(n)=2^nt_{n+1}$ for all $n\ge0$, and $a(n+1)/a(n)=2t_{n+2}/t_{n+1}\to2$.\end{theorem}
\begin{proof} An integer $m$ with $\Omega(m)=n$ is uniquely $m=2^{n-j}u$ with $u$ odd and $\Omega(u)=j\le n$, so
$m\mapsto m/2^n=u/2^j$ is an increasing bijection from $\{m:\Omega(m)=n\}$ onto $V_n=\{u/2^{\Omega(u)}:\Omega(u)\le n\}$.
If $\Omega(u)=j$ then $u\ge3^j$ and $u/2^j\ge(3/2)^j$, so the elements of $V\setminus V_n$ exceed $(3/2)^n$. On the
other hand $t_{n+1}\le p_{n+2}/2$, since $1$ and the $p/2$ ($p$ an odd prime) are in $V$; for $n\ge10$,
$p_{n+2}/2<(3/2)^n$, so the $n+1$ smallest elements of $V$ lie in $V_n$ and $a(n)=2^nt_{n+1}$. For $n<10$ (indeed
for the 30 known terms) this is checked directly. $V$ is locally finite, since $u/2^j\le y$ forces
$(3/2)^j\le y$ and $u\le2^jy$; so $t_k\to\infty$. Given $\varepsilon>0$, the prime number theorem gives, for $y$
large, a prime in $(2y,2y(1+\varepsilon)]$, that is, an element of $V$ in $(y,y(1+\varepsilon)]$. With $y=t_k$ we get
$t_k\le t_{k+1}\le(1+\varepsilon)t_k$ for large $k$, hence $t_{k+1}/t_k\to1$. \end{proof}

\section{Excursions in a cone: Kotesovec's exponent (\oeis{A361190})}\label{s:A361190}
$a(n)$ counts the excursions of length $4n$ in the wedge $0\le y\le x$ with steps $(1,1)$, $(1,-1)$, $(-1,0)$.
Kotesovec observed $a(n)\sim c\,64^n/n^\alpha$ and conjectured $\alpha=1+\pi/\arctan\sqrt2\approx4.288536$.
\begin{proposition} The conjecture holds, by the theorem of Denisov and Wachtel~\cite{DW}.\end{proposition}
\begin{proof}[Argument] The map $(u,v)=(x-y,y)$ sends the wedge to the quarter plane and the steps to
$(0,1)$, $(2,-1)$, $(-1,0)$. The exponential weights $e^{a\,dx+b\,dy}$ cancel the drift for $b=0$ and
$e^{2a}=1/2$; the probabilities become $(1/4,1/4,1/2)$ and the growth rate is $2\sqrt2$ per step, that is
$64^n$ for $4n$ steps. The weight of an excursion is $1$, so $a(n)=64^n\,P(\text{excursion of the centred
walk})$. The centred walk has $\mathrm{Var}\,u=3/2$, $\mathrm{Var}\,v=1/2$, $\mathrm{Cov}=-1/2$, correlation
$-1/\sqrt3$; after decorrelation the quarter plane becomes a cone of angle
$\arccos(1/\sqrt3)=\arctan\sqrt2$. For a zero-drift walk with bounded steps, Denisov and Wachtel give
$P(S_N=0,\ \tau>N)\asymp N^{-1-\pi/\theta}$ along the lattice of returns ($N\in4\mathbb{Z}$ here), hence
$\alpha=1+\pi/\arctan\sqrt2$. \end{proof}
As $\arctan(\sqrt2)/\pi=\arccos(-1/3)/(2\pi)$ is irrational, $\alpha$ is irrational and the generating function
is not D-finite (the argument of Bostan--Raschel--Salvy~\cite{BRS}). This is an application of a published
theorem, not a self-contained proof; the hypotheses (strong aperiodicity on the lattice generated by the
steps, non-degenerate covariance) are standard for bounded steps but were not re-checked line by line.
Numerically, the local slope of $\log(a(n)/64^n)$ is $4.236$, $4.271$, $4.278$, $4.290$ at $n=100,300,500,599$.

\section{The Thue--Morse convolution inverse (\oeis{A225132})}\label{s:A225132}
The data of \oeis{A225132} are the coefficients of $1/h(x)$, where $h(x)=\sum_{n\ge0}t(n+1)x^n=g(x)/x$ and
$g$ is the generating function of \oeis{A010060} ($t(0)=0$); the comment ``$1/(1+g(x))$'' does not reproduce
the data. The conjecture is $a(n+1)/a(n)\to-1.405177106052\ldots$
\begin{proposition} This is true: $a(n+1)/a(n)\to1/x_0$, where $x_0=-0.7116540653082362\ldots$ is the only zero
of $h$ in $|x|<0.76$.\end{proposition}
\begin{proof}[Proof (computer-assisted, \texttt{A225132/proof\_check.py})] The truncation $h_N$ with $N=600$
terms differs from $h$ by at most $r^{N+1}/(1-r)$ on $|x|\le r$. A sign change of $h$ is certified on an
interval of width $2\cdot10^{-12}$ around $x_0$, and $h'(x_0)\approx1.8654\ne0$. A certified argument principle
(40\,000 points on each circle; the minimum of $|h_N|$ exceeds the sum of the truncation error and of the
variation between points, bounded with $|h'|\le1/(1-r)^2$) shows that $h$ has no zero in $|x|<0.70$ and exactly
one in $|x|<0.76$, which is $x_0$. Hence $1/h-c/(x-x_0)$, with $c=1/h'(x_0)$, is analytic in $|x|<0.76$, its
coefficients are $O(0.75^{-n})$, while those of $c/(x-x_0)$ are $-c\,x_0^{-n-1}$, of modulus $\asymp1.405^n$. So
$a(n)\sim-x_0^{-n-1}/h'(x_0)$ and $a(n+1)/a(n)\to1/x_0=-1.405177106052044\ldots$ \end{proof}

\section{Binary partitions: Cloitre's constant (\oeis{A033485})}\label{s:A033485}
Let $a(1)=1$ and $a(n)=a(n-1)+a(\lfloor n/2\rfloor)$; for $n\ge1$, $a(n)$ is half the number of binary partitions
of $2n$ (\oeis{A000123}). Cloitre (2003) conjectured that $a(2n)/a(n)\cdot\log(n)/n$ tends to $c=1.64\ldots$
\begin{theorem}\label{t:A033485} Let $\beta(n)=n\,a(\lfloor n/2\rfloor)/a(n)$. Then
$\beta(n)=\log_2n+O((\log\log n)^2)$. Consequently
\[ \frac{a(2n)}{a(n)}\cdot\frac{\log n}{n}=\frac{2\log n}{\beta(2n)}\longrightarrow2\log2=1.38629\ldots, \]
so the limit exists, and Cloitre's value $1.64\ldots$ is wrong.\end{theorem}
The value $1.64$ came from $n\le10^4$: the quantity has a maximum of about $1.659$ near $n=150$, crosses $1.64$
near $n=10^4$, and is $1.566$ at $n=1.5\cdot10^8$, still decreasing.

\begin{proof}[Proof (computer-assisted)] \emph{Exact identities.} $a(n-1)/a(n)=1-\beta(n)/n$, and
$0<\beta(n)/n\le1/2$ since $a(n)\ge2a(\lfloor n/2\rfloor)$. Each $k<N$ appears twice and $N$ once in
$a(2N)=1+\sum_{j=2}^{2N}a(\lfloor j/2\rfloor)$, so $a(2N)=2S(N)-a(N)+1$ and $a(2N+1)=a(2N)+a(N)$, with
$S(N)=a(1)+\dots+a(N)$. Put $\sigma(N)=S(N)/a(N)=\sum_{k\le N}P(k)$ with $P(k)=\prod_{j=k+1}^N(1-\beta(j)/j)$, and
$\rho=N/\sigma(N)$. Then
\[ \beta(2N)=\frac{N}{\sigma-\frac12+\frac1{2a(N)}}\in\Bigl[\rho,\frac{\rho}{1-\rho/(2N)}\Bigr],\qquad
   \beta(2N+1)=\frac{N+\frac12}{\sigma+\frac1{2a(N)}}\in\Bigl[\frac{\rho}{1+\rho/(2N)},\rho\bigl(1+\tfrac1{2N}\bigr)\Bigr]. \]

\emph{Local step.} Let $N\ge2^{28}$, $L=\log_2N$, and assume $|\beta(j)-\log_2j|\le D$ for $j\in(N/2,N]$, with
$D\le L/2$. For $\varepsilon\in(0,1/2)$ and $k_1=\lceil(1-\varepsilon)N\rceil$, every $j>k\ge k_1$ has
$\beta(j)\ge b^*=L+\log_2(1-\varepsilon)-D$, so $P(k)\le((k+1)/(N+1))^{b^*}$, and $P(k)\le P(k_1)$ for $k<k_1$; hence
$\sigma\le(N+1)/(b^*+1)+1+N(1-\varepsilon+2/N)^{b^*}$. In the other direction, for $k\ge N/2$ and $j\in(k,N]$,
$\beta(j)\le B=L+D$, and $1-x\ge e^{-x-x^2}$ gives $P(k)\ge(k/N)^Be^{-2B^2/N}$, so
$\sigma\ge e^{-2B^2/N}N(1-2^{-B-1})/(B+1)$. With the identities, for $n\in\{2N,2N+1\}$:
$|\beta(n)-\log_2n|\le D+\delta(L,D)$, with an explicit $\delta$ (function \texttt{delta} of
\texttt{A033485/bootstrap.py}, with $\varepsilon=t\ln b_0/b_0$, $t$ optimised, $b_0=L-D$); asymptotically
$\delta\approx(3/\ln2)\ln b_0/b_0+1/b_0$, and $\delta\le9.5\ln L/L$ when $D\le L/2$ and $L\ge10^4$.

\emph{Induction on dyadic blocks.} Let $E_k=\max|\beta(j)-\log_2j|$ over $(2^k,2^{k+1}]$. Since an $n$ in
$(2^{k+1},2^{k+2}]$ is $2N$ or $2N+1$ with $(N/2,N]\subset(2^{k-1},2^{k+1}]$,
$E_{k+1}\le\max(E_k,E_{k-1})+\delta(k,\max(E_k,E_{k-1}))$ while the hypotheses hold. The base is computed in double
precision: $E_{27}=4.1068$, $E_{28}=4.1618$. Iterating the bound from $k=28$ to $10^4$ gives $D_k/k\le0.355$
(maximum near $k=100$, then decreasing to $0.02$ at $k=10^4$), so $D\le L/2$ always holds; beyond $10^4$,
$\delta\le9.5\ln k/k$ gives $D_k\le D_{10^4}+4.75(\ln^2k-\ln^210^4)\ll k/2$. Hence $E_k=O(\log^2k)$ and
$\beta(n)/\log_2n\to1$. \end{proof}
\begin{remark}[A shorter route, and prior work] The limit also follows from de Bruijn's asymptotic expansion
for binary partitions (1948), recorded by Flajolet in \oeis{A000123}: with $L_2=\log2$,
$\log a(n)=\frac1{2L_2}\log^2\frac n{\log n}+\bigl(\frac12+\frac{1+\log L_2}{L_2}\bigr)\log n-\bigl(1+\frac{\log L_2}{L_2}\bigr)\log\log n
+\Phi\bigl(\log\frac n{\log n}/L_2\bigr)+o(1)$ with $\Phi$ of period $1$. Taking the difference between $2n$ and $n$, the
periodic terms cancel up to $o(1)$ and $\log a(2n)-\log a(n)=\log n-\log\log n+\log(2\log2)+o(1)$, which gives the
limit $2\log2$. So Cloitre's value is refuted by a classical theorem, and the argument below is an independent,
self-contained check.\end{remark}
\begin{remark} The self-contained proof is not fully certified: the base values $E_{27},E_{28}$ are computed in floating point
(relative error about $10^{-7}$, against a margin of $10^{-5}$); the worst case $L=k$ inside a block is checked on a
grid ($k=28,50,200,1000$); and the smallest margin $k/2-D_k$ in the iteration is $1$. A certified version would
redo the base in interval arithmetic. The result is consistent with de Bruijn's asymptotics for binary partitions
(1948), which we did not use. The second conjecture of the same line ($a(n)/A(n)$ bounded) is not treated.
\end{remark}

\section{Products of Wythoff rows (\oeis{A022344})}\label{s:A022344}
For a positive Fibonacci sequence $T$ let $Q(T)=|T(i)^2-T(i-1)T(i+1)|$, independent of $i$; \oeis{A022344}$(n)$ is
$Q$ of row $n$ of the Wythoff array. Ramsey (2007) conjectured: if $F_1$, $F_2$ lie in rows $n$, $m$, then for fixed
$j$ the sequence $F_3(i)=F_1(i)F_2(j+1)-F_1(i+1)F_2(j)$ lies (eventually) in a row $p$ with $a(p)=a(n)a(m)$.
\begin{proposition} This is true.\end{proposition}
\begin{proof} $F_3$ is a fixed linear combination of $F_1(i)$ and $F_1(i+1)$, hence a Fibonacci sequence in $i$; it
is also one in $j$. With $x=G(i-1)$, $y=G(i)$, $G(i)^2-G(i-1)G(i+1)=y^2-xy-x^2$ is the norm form of
$\mathbb{Z}[\varphi]$, which changes only by a sign when $i$ moves. So $|Q(F_3)|$ does not depend on $i$ or $j$, and
the identity $|Q(F_3)|=|Q(F_1)|\,|Q(F_2)|$, a polynomial identity in the four initial values, is checked
symbolically (\texttt{A022344/check.py}). Since $y^2-xy-x^2$ has no non-trivial integer zero, $F_3\not\equiv0$, and
$F_3(i)\approx F_1(i)(F_2(j+1)-\varphi F_2(j))$ with a non-zero factor, so $F_3$ has eventually constant sign. By
Morrison's theorem (1980), every positive Fibonacci sequence eventually coincides with a row $p$ of the Wythoff
array, and $a(p)=|Q(F_3)|=a(n)a(m)$. \end{proof}
The map $(n,m)\mapsto p$ is \oeis{A357097} (Section~\ref{s:A357097}); the entry notes that
$a(\oeis{A357097}(x,y))=a(x)a(y)$ (Munn, 2025), which is the core of the conjecture. Contribution: modest.
